\documentclass[8.2pt]{article}
\usepackage[utf8]{inputenc}
\usepackage[top=0.35in, bottom=0.35in, left=0.4in, right=0.4in]{geometry}
\usepackage{enumitem}
\usepackage[hidelinks]{hyperref}
\usepackage{parskip}
\setlength{\parindent}{0pt}

\renewcommand{\familydefault}{\sfdefault}
\newcommand{\sect}[1]{\vspace{1pt}\noindent\textbf{\large #1}\vspace{1pt}\hrule\vspace{3pt}}

\begin{document}

% ---------- Header ----------
\begin{center}
  {\Large \textbf{Thirulok Sundar Mohan Rasu}}\\[4pt]
  Ann Arbor, MI 48105 $\,|\,$ +1 (248) 704-2911 $\,|\,$
  \href{mailto:thirulok@umich.edu}{thirulok@umich.edu} $\,|\,$
  \href{https://thiruloksundar.github.io}{thiruloksundar.github.io} $\,|\,$
  \href{https://github.com/Thiruloksundar}{github.com/Thiruloksundar} $\,|\,$
  \href{https://www.linkedin.com/in/thirulok-sundar-mohanrasu/}{linkedin.com/in/thirulok-sundar-mohanrasu}
\end{center}
\vspace{-3mm}

% ---------- Education ----------
\sect{Education}
\vspace{-2mm}
\textbf{University of Michigan -- Ann Arbor} \hfill Aug 2025 -- Dec 2026\\
M.S., Electrical and Computer Engineering \hfill GPA: 4.0 / 4.0\\
Relevant courses: \emph{Matrix Methods for ML \& Signal Processing (ECE 551), Probability \& Random Processes (ECE 501)}\\[4pt]
\textbf{Indian Institute of Technology (ISM) Dhanbad} \hfill Dec 2021 -- May 2025\\
B.Tech., Electrical Engineering. Member: \emph{CyberLabs (Machine Learning)} \hfill GPA: 8.30 / 10

% ---------- Skills ----------
\sect{Skills}
\vspace{-2mm}
\textbf{Languages:} Python, SQL, Julia, MATLAB\\
\textbf{Data Science \& ML:} Pandas, NumPy, PySpark, \textbf{Machine Learning}, Data Ingestion \& Consolidation, Feature Engineering, Data Visualization (Matplotlib, Seaborn), JSON/NoSQL, Time-Series Analysis\\
\textbf{Tools:} PyTorch, AWS S3, Scikit-learn, SciPy, HuggingFace, GitHub, Linux

% ---------- Experience ----------
\sect{Experience}
\vspace{-2mm}

\textbf{Electric Vehicle Center -- University of Michigan} \hfill Oct 2025 -- Present\\
\emph{Graduate Research Assistant (Prof.\ Ziyou Song)}\\
\vspace{-6mm}
\begin{itemize}[leftmargin=*,noitemsep,topsep=1pt]
  \item Built \textbf{data ingestion and processing pipelines} in \textbf{Python/Pandas} to collect, clean, and consolidate raw \textbf{electrochemical time-series datasets} (117 EIS spectra across 6 cells, multiple SOC/SOH levels); conducted \textbf{feature sensitivity (OAT) analysis} to identify most impactful parameters (SEI thickness, CPE exponent) driving impedance response from a 6-dimensional feature space
  \item Performed \textbf{feature engineering and selection} on early-cycle battery data --- identified discharge capacity (\texttt{qdis}), cumulative Ah throughput, C-rates, depth of discharge, and DVA features as key predictors via \textbf{PCA-based dimensionality reduction}; used these as inputs to an \textbf{ML surrogate model} (attention-enhanced Seq2Seq) achieving \textbf{2--5\% MAPE} on capacity-fade trajectories across 225 cells
  \item Developed \textbf{data-driven optimization framework} in Python to analyze accelerated test protocols; processed real-world driving profiles into structured \textbf{battery power demand datasets}, enabling systematic analysis of degradation mechanism diversity (LAM-NE, LAM-PE, SEI, LLI) across experimental conditions
\end{itemize}

\vspace{-2mm}
\textbf{Carnegie Mellon University} \hfill Jan 2024 -- Sep 2025\\
\emph{Research Intern (Dr.\ Arun Balajee Vasudevan)}\\
\vspace{-6mm}
\begin{itemize}[leftmargin=*,noitemsep,topsep=1pt]
  \item Engineered \textbf{end-to-end data generation and labeling pipeline} in Python to curate a \textbf{1M+ sample dataset} of real/synthetic audio pairs; data stored in structured \textbf{JSON (NoSQL) format} with transcript, label, and metadata fields --- enabling scalable ingestion and querying
  \item Fine-tuned \textbf{CLAP} audio-language model on curated dataset for fake audio detection achieving \textbf{90\% accuracy}; built automated evaluation pipelines for benchmarking across multiple TTS systems \emph{\textbf{\href{https://drive.google.com/file/d/1hdSPxQQA-PKayYQPPESZb11ofm7XrQRX/view?usp=sharing}{(paper)}}}
\end{itemize}

\vspace{-2mm}
\textbf{University of Central Florida -- CRCV Lab} \hfill Feb 2024 -- Dec 2024\\
\emph{Research Intern (Prof.\ Yogesh Rawat)}\\
\vspace{-6mm}
\begin{itemize}[leftmargin=*,noitemsep,topsep=1pt]
  \item Created and \textbf{analyzed a large annotated video dataset} with structured \textbf{JSON entries} capturing actor, action, object, and environment labels per clip; designed fine-grained evaluation metrics and \textbf{QA processes} for activity understanding in Video-LLMs
  \item Defined \textbf{data schema} and annotation taxonomy to standardize labeling across annotators; contributed to \textbf{data consolidation} pipeline integrating multi-source video collections into a unified evaluation benchmark
\end{itemize}

\vspace{-2mm}
\textbf{Mowito} \hfill Mar 2023 -- Jul 2023\\
\emph{Software Intern}\\
\vspace{-6mm}
\begin{itemize}[leftmargin=*,noitemsep,topsep=1pt]
  \item Processed and labeled robot vision data stored in \textbf{AWS S3 buckets}; built modular \textbf{data annotation pipeline} supporting multiple formats (instance segmentation, picking, crate detection) with visualization tools; integrated into production pipelines for warehouse automation
  \item Implemented \textbf{perceptual hashing} for large-scale \textbf{dataset deduplication}, reducing dataset size by \textbf{30\%}; developed \textbf{order-fulfillment optimization} using \textbf{Leiden clustering} to group crates and minimize completion time across high-density SKU environments
\end{itemize}

% ---------- Projects ----------
\sect{Selected Projects}
\vspace{-2mm}

\textbf{ML Pipeline for Battery Aging Prediction} \hfill 2025\\
Built full \textbf{data processing and ML pipeline} in Python: ingested raw NMC cycling data, performed \textbf{feature engineering} (discharge capacity, Ah throughput, C-rate, DoD), applied \textbf{PCA-based feature selection}, and trained attention-enhanced \textbf{Seq2Seq model} for \textbf{time-series forecasting}; visualized degradation trajectories and error distributions across 225 cells using Matplotlib/Seaborn.

\textbf{Multi-view 3D Scene Reconstruction Analysis} \hfill 2025\\
Collected and annotated a structured scene dataset; performed systematic \textbf{data-driven analysis} of reconstruction failure cases across varying lighting, occlusion, and viewpoint conditions; documented findings with quantitative metrics and visualizations.

\end{document}
